\begin{slide}
\begin{center}\Title{Concluding remarks}\end{center}

\begin{itemize}

\IncPart{1}{\item[]
Observations 1 and 2 imply that the construct $\IRef{h}T$
denotes the application of some $\Uniform{h\p}$ to $T$,
where $h\p$ depends on context of $\IRef{h}T$. 
}

\IncPart{2}{\item[]
An \stress{algorithm} exists to check if a path $r \in T$ has the form
$\IS\LabP{p}\LabA\LabP{b}\LabL\LabP{q}\\\LabP{(n+1)}$
in which $b$ and $q$ meet the requirements for ef-reduction.
}

\IncPart{3}{\item[]
We can give a variant of the ef-reduction grounded in the assumption that
$\IRef{h}T$ denotes the context independent application of $\Uniform{h}$ to $T$.
}

\IncPart{3}{\item[]
Under the assumptions of ef-reduction,
$T \EFpStep{r} T\Subs{\IS\LabP{r\p}\Subtree{T}(\IS\LabP{p}\LabS)}{r}$\\
where $r\p \bydef \IS\LabP{p}\LabA\LabP{\Clear{b}}\LabL\LabP{q}\LabP{(\Length{b}+n+1)}$
and $\Clear{b}$ sets to $0$ the labels $h$ of $b$. 
}

\IncPart{4}{\item[]
To \stress{apply} the delayed lifts in a path $\IS\LabP{p}\LabD{n} \in T$
we remove them from $p$ and we apply to $n$ an update function (unfold)
depending on $p$.
}

\IncPart{5}{\item[]
We can express other reductions in this system as well.
For instance variants of non-expanding focused $\beta$-reduction and $\zeta$-reduction.
}

\IncPart{6}{\item[]
We are verifying that all these reductions have desired properties,
e.g., \stress{confluence}.
Such results will appear in a forthcoming article.
}

\end{itemize}

\end{slide}
